% SOURCE_AUTHORITY: sga5_fr_workpass.tex, lines 5311--5345 (LF-normalized unit).
% SOURCE_SHA256: 791F4EFFC5E02832D5D77ED03518C8156D6F07E4C8238B03545DB93D883FBB28
La primera afirmación de 5.10.9 es trivial. Demostremos la segunda. Se puede suponer que $M$ tiene componentes inyectivas y que $E$ tiene componentes de la forma $B\otimes_KL$, con $L$ plano. Se tiene entonces
\[
\uRHom_B(E,B\lten_KM)\lten_BE=\Hom_B^\bullet(E,B\otimes_KM)\otimes_BE,
\]
\[
\uRHom_A(E,A\lten_KM)\lten_AE=\Hom_A^\bullet(E,A\otimes_KM)\otimes_AE,
\]
y las flechas de 5.10.10 están definidas en el nivel de los complejos. Quedamos, pues, reducidos a comprobar que, si $L$, $M$ son $K$-módulos y $E$ es un $B$-módulo izquierdo, se tiene un cuadrado conmutativo
\[
\begin{tikzcd}[column sep=large,row sep=large]
|[alias=star]| {(*)} & \uHom_B(B\otimes_KL,B\otimes_KM)\otimes_B(B\otimes_KL) \arrow[r,"5.3.3"] \arrow[d,"(1)"'] & M\otimes_KB_\natural \arrow[d,"M\otimes_K\Tr_{B/A}"]\\
& \uHom_A(B\otimes_KL,A\otimes_KM)\otimes_A(B\otimes_KL) \arrow[r,"5.3.3"'] & M\otimes_KA_\natural,
\end{tikzcd}
\]
donde la flecha (1) se define enviando $B\otimes_KM$ a $\uHom_A(B,A)\otimes_AB\otimes_KM$ mediante 5.10.1 y componiendo con un isomorfismo del tipo 5.10.3 y un isomorfismo de adjunción. Como $A$ es plano y de presentación finita como $K$-módulo, las flechas ``producto tensorial''
\[
\text{(2)}\quad B\otimes_K\uHom_K(L,M)=\uHom_B(B,B)\otimes_K\uHom_K(L,M)\to\uHom_B(B\otimes_KL,B\otimes_KM),
\]
\[
\text{(3)}\quad \uHom_A(B,A)\otimes_K\uHom_K(L,M)\to\uHom_A(B\otimes_KL,A\otimes_KM)
\]
son isomorfismos, mediante los cuales (1) se identifica con el producto tensorial (sobre $K$) de 5.10.1 por $\uHom_K(L,M)\otimes_KL$. Por otra parte, la flecha horizontal superior (resp. inferior) de (*) se identifica, mediante (2) (resp. (3)), con el producto tensorial (sobre $K$) de $\Tr_B:B\to B_\natural$ (resp. $\mathrm{ev}:\uHom_A(B,A)\otimes_AB\to A_\natural$) por la flecha de evaluación $\uHom_K(L,M)\otimes_KL\to M$. La conmutatividad de $(*)$ se deduce, por tanto, de la de 5.10.2, lo que termina la demostración de 5.10.9.

\begin{statement}{Corolario 5.10.11}
Bajo las hipótesis de 5.10.9, para $E\in\ob D(B)_{\mathrm{parf}}$, $M\in\ob D^b(K)$, se tiene un cuadrado conmutativo de $D(K)$
\begin{equation}\tag{5.10.12}
\begin{tikzcd}[column sep=large,row sep=large]
\uRHom_B(E,M\lten_KE) \arrow[r,"\Tr_B"] \arrow[d] & M\lten_KB_\natural \arrow[d,"M\lten\Tr_{B/A}"]\\
\uRHom_A(E,M\lten_KE) \arrow[r,"\Tr_A"'] & M\lten_KA_\natural,
\end{tikzcd}
\end{equation}
donde la flecha vertical de la izquierda es la restricción de escalares.
\end{statement}

En efecto, basta conjugar 5.10.6 y 5.10.9.
